\documentclass[../../main.tex]{subfiles}
\graphicspath{{../../image/}{../../figure/}}
\begin{document}

\section{动力论表述}\label{section:ch06-动力论表述}

本节给出\eqref{align:ch06-6-1-1}、\eqref{align:ch06-6-1-2}可容许弱解的另一种等价刻画，下文将看到它由动力论启发。

早已注意到，对满足\eqref{align:ch06-6-2-2}的\eqref{align:ch06-6-1-1}任一解，熵产生是非正测度。特别地，若 $u$ 是\cref{definition:ch06-6-2-1}意义下\eqref{align:ch06-6-1-1}、\eqref{align:ch06-6-1-2}的可容许解，则对任一 $v\in(-\infty,\infty)$，
\begin{align}
\partial_{t}\{|u-v|-|v|\}+\operatorname{div}\{\operatorname{sgn}(u-v)[G(u)-G(v)]-\operatorname{sgn}v\,G(v)\}=-2\nu_{v},
\label{align:ch06-6-7-1}
\end{align}
其中 $\nu_{v}$ 是 $\mathbb{R}^{m}\times\mathbb{R}^{+}$ 上的非负测度。当 $|v|>\sup|u_{0}|=\sup|u|$ 时，$2\nu_{v}=|\partial_{t}u+\operatorname{div}G(u)|=0$。

把 $\{\nu_{v}\}$ 实现为 $\mathbb{R}\times\mathbb{R}^{m}\times\mathbb{R}^{+}$ 上的非负测度 $\nu$，在分布意义下关于 $v$ 微分\eqref{align:ch06-6-7-1}，得
\begin{align}
\partial_{t}\chi(v;u)+\sum_{\alpha=1}^{m}G_{\alpha}^{\prime}(v)\partial_{\alpha}\chi(v;u)=\partial_{v}\nu,
\label{align:ch06-6-7-2}
\end{align}
其中 $\chi$ 表示函数
\begin{align}
\chi(v;u)=
\begin{cases}
1, & 0<v<u,\\
-1, & u<v<0,\\
0, & \text{其他}.
\end{cases}
\label{align:ch06-6-7-3}
\end{align}

任一熵--熵流对 $(\eta,Q)$ 的熵产生都易于用 $\nu$ 表示。事实上，以 $\eta^{\prime}(v)$ 乘\eqref{align:ch06-6-7-2}并关于 $v$ 在 $(-\infty,\infty)$ 上积分。回顾\eqref{align:ch06-6-2-1}并分部积分，得
\begin{align}
\partial_{t}\int_{-\infty}^{\infty}\eta^{\prime}(v)\chi(v;u)\,\mathrm{d}v+\operatorname{div}\int_{-\infty}^{\infty}Q^{\prime}(v)\chi(v;u)\,\mathrm{d}v=-\int_{-\infty}^{\infty}\eta^{\prime\prime}(v)\,\mathrm{d}\nu(v;\cdot,\cdot).
\label{align:ch06-6-7-4}
\end{align}
容易验证，若 $p(v)$ 是任一 $C^{1}$ 函数，则
\begin{align}
\int_{-\infty}^{\infty}p^{\prime}(v)\chi(v;u)\,\mathrm{d}v=p(u)-p(0),
\label{align:ch06-6-7-5}
\end{align}
故\eqref{align:ch06-6-7-4}给出
\begin{align}
\partial_{t}\eta(u)+\operatorname{div}Q(u)=-\int_{-\infty}^{\infty}\eta^{\prime\prime}(v)\,\mathrm{d}\nu(v;\cdot,\cdot).
\label{align:ch06-6-7-6}
\end{align}
特别地，当 $\eta(u)$ 为凸时，\eqref{align:ch06-6-7-6}右端非正。此外，取 $\eta(u)=\tfrac{1}{2}u^{2}$ 并关于 $(x,t)$ 在 $\mathbb{R}^{m}\times[0,\infty)$ 上积分，得
\begin{align}
\int_{0}^{\infty}\int_{\mathbb{R}^{m}}\int_{-\infty}^{\infty}\mathrm{d}\nu(v;x,t)\,\mathrm{d}x\,\mathrm{d}t\leqslant\frac{1}{2}\int_{\mathbb{R}^{m}}u_{0}^{2}(x)\,\mathrm{d}x.
\label{align:ch06-6-7-7}
\end{align}

值得注意的是，\eqref{align:ch06-6-7-2}完全刻画了\eqref{align:ch06-6-1-1}的可容许弱解，如下述定理所示。

\begin{theorem}\label{theorem:ch06-6-7-1}
$\mathbb{R}^{m}\times[0,\infty)$ 上的有界可测函数 $u(x,t)$ 是\eqref{align:ch06-6-1-1}、\eqref{align:ch06-6-1-2}的可容许解，当且仅当由\eqref{align:ch06-6-7-3}定义的函数 $\chi(v;u(x,t))$ 在 $\mathbb{R}\times\mathbb{R}^{m}\times[0,\infty)$ 上对某非负测度 $\nu$ 满足\eqref{align:ch06-6-7-2}，并满足初始条件
\begin{align}
\chi(v;u(x,0))=\chi(v;u_{0}(x)),\qquad v\in(-\infty,\infty),\ x\in\mathbb{R}^{m}.
\label{align:ch06-6-7-8}
\end{align}
\end{theorem}

\begin{proof}
方程\eqref{align:ch06-6-7-2}容许解 $\chi(\cdot;u(\cdot,t))\in C^{0}([0,\infty);L^{1}(\mathbb{R}\times\mathbb{R}^{m}))$，因而初始条件\eqref{align:ch06-6-1-2}在 $L^{1}(\mathbb{R}^{m})$ 中强成立。余下只需证明对每个凸熵--熵流对 $(\eta,Q)$（$\eta$ 凸），\eqref{align:ch06-6-2-2}成立。由于 $u$ 有界，只需对线性增长的熵（即 $|\eta^{\prime}(u)|$ 在 $(-\infty,\infty)$ 上有界）建立\eqref{align:ch06-6-2-2}。

从\eqref{align:ch06-6-7-2}出发，可如前证明\eqref{align:ch06-6-7-6}成立，但仅对那些 $\eta^{\prime}(v)$ 在 $|v|$ 充分大时为零的函数 $\eta(v)$ 成立（以便分部积分；因不再已知 $\nu$ 在 $|v|>\sup|u_{0}|$ 时为零）。

固定任一具有线性增长的凸函数 $\eta$，对 $k=1,2,\cdots$ 令 $\eta_{k}(v)=\eta(v)\phi(v/k)$，其中 $\phi$ 是 $(-\infty,\infty)$ 上光滑偶函数，$|v|\leqslant 1$ 时 $\phi(v)=1$，$|v|\geqslant 2$ 时 $\phi(v)=0$，$v\in(1,2)$ 时 $\phi^{\prime}(v)<0$。于是
\begin{align}
\partial_{t}\eta_{k}(u)+\operatorname{div}Q_{k}(u)&=-\int_{-\infty}^{\infty}\eta_{k}^{\prime\prime}(v)\,\mathrm{d}\nu(v;\cdot,\cdot)\notag\\
&=-\int_{-\infty}^{\infty}\Bigl{[}\eta^{\prime\prime}(v)\phi\Bigl{(}\frac{v}{k}\Bigr{)}+\frac{2}{k}\eta^{\prime}(v)\phi^{\prime}\Bigl{(}\frac{v}{k}\Bigr{)}+\frac{1}{k^{2}}\eta(v)\phi^{\prime\prime}\Bigl{(}\frac{v}{k}\Bigr{)}\Bigr{]}\,\mathrm{d}\nu(v;\cdot,\cdot).
\label{align:ch06-6-7-9}
\end{align}
当 $k$ 充分大时，在解的值域上 $\eta_{k}(u)=\eta(u)$、$Q_{k}(u)=Q(u)$。此外，$\eta^{\prime\prime}(v)\phi(v/k)$ 当 $k\to\infty$ 时单调趋于 $\eta^{\prime\prime}(v)$。最后显然有 $\eta^{\prime}(v)\phi^{\prime}(v/k)=O(1)$、$\eta(v)\phi^{\prime\prime}(v/k)=O(k)$（$k\to\infty$）。于是在\eqref{align:ch06-6-7-9}中令 $k\to\infty$，即得\eqref{align:ch06-6-7-6}，从而得\eqref{align:ch06-6-2-2}。证明完成。
\end{proof}

动力论表述\eqref{align:ch06-6-7-2}可作为\eqref{align:ch06-6-1-1}可容许弱解的另一种等价定义，为发现这些解的性质提供了强有力的工具。特别地，即使在初值仅属于 $L^{1}(\mathbb{R}^{m})$ 而不必属于 $L^{\infty}(\mathbb{R}^{m})$ 的更一般假设下，也可获得 L$^{1}$ 压缩性质\eqref{align:ch06-6-2-9}的另一直接证明；文献见\cref{section:ch06-注记}。


至此，我们一直把非线性视为在光滑初值的解中诱发间断发展的因素。然而结果表明，非线性也可能扮演相反的角色：磨光粗糙初值的解。在本书后续中将遇到此类行为的各种表现。动力论表述为洞察标量守恒律中非线性的紧化与磨光效应提供了宝贵视角。这将在下节\cref{section:ch06-Linfty解精细结构}中变得明显，但也见于下述正则性定理，其困难而技术性的证明见\cref{section:ch06-Linfty解精细结构}所引文献。


\begin{theorem}\label{theorem:ch06-6-7-2}
假设存在 $r\in(0,1]$ 与 $C\geqslant 0$，使对所有 $\delta\in(0,1)$、$p\in\mathbb{R}$、$P\in\mathbb{R}^{m}$（$p^{2}+|P|^{2}=1$），
\begin{align}
\operatorname{meas}\{v:|v|\leqslant\|u_{0}\|_{L^{\infty}},\ |p+G^{\prime}(v)P|\leqslant\delta\}\leqslant C\delta^{r}.
\label{align:ch06-6-7-10}
\end{align}
则\eqref{align:ch06-6-1-1}、\eqref{align:ch06-6-1-2}的可容许弱解 $u$ 满足
\begin{align}
u(\cdot,t)\in C^{0}\bigl{(}(0,\infty);W_{loc}^{s,1}(\mathbb{R}^{m})\bigr{)},
\label{align:ch06-6-7-11}
\end{align}
对任意 $s\in(0,\frac{r}{2r+1})$。
\end{theorem}

正是条件\eqref{align:ch06-6-7-10}编码了 $G$ 中负责正则化效应的非线性方面。例如，当 $G$ 为线性时，对任何 $r$，\eqref{align:ch06-6-7-10}不成立；但若 $G_{\alpha}$ 为一致凸函数（$G_{\alpha}^{\prime\prime}(u)>0$，$\alpha=1,\cdots,m$），则取 $r=1$ 时\eqref{align:ch06-6-7-10}成立。

本节末尾讨论标量守恒律的动力论表述\eqref{align:ch06-6-7-2}如何由物质动力论启发。正如在\cref{chapter:ch03-双曲平衡律方程组}中所见，在经典气体动力论中，气体在点 $x$、时刻 $t$ 的状态由分子速度 $\xi$ 的分子密度函数 $f(\xi,x,t)$ 描述。

$f$ 的演化由 Boltzmann 方程\eqref{align:ch03-Boltzmann}支配，它监测分子速度分布因输运与碰撞而发生的变化。

动力论途径与连续介质途径之间的联系，是通过把强度量（如密度、速度、压强、温度等）与分子密度函数 $f$ 的适当矩等同起来，再证明这些场满足连续介质物理的平衡律而建立的。因而原则上，可把场视为某下层动力论模型中分子密度的矩，从而构造平衡律方程组的解；该模型的密度函数的零阶矩满足标量守恒律\eqref{align:ch06-6-1-1}。

在模型中，“速度” $v$ 取标量值，点 $x$、时刻 $t$ 的“分子密度” $f(v;x,t)$ 允许取正负值。于是 $u$ 由 $f$ 经
\begin{align}
u(x,t)=\int_{-\infty}^{\infty}f(v;x,t)\,\mathrm{d}v
\label{align:ch06-6-7-12}
\end{align}
得到。而 $f$ 满足输运方程
\begin{align}
\partial_{t}f(v;x,t)+\sum_{\alpha=1}^{m}G_{\alpha}^{\prime}(v)\partial_{\alpha}f(v;x,t)=\frac{1}{\mu}[\chi(v;u(x,t))-f(v;x,t)],
\label{align:ch06-6-7-13}
\end{align}
其中 $\mu$ 是小正参数，$\chi(v;u)$ 是\eqref{align:ch06-6-7-3}所定义的函数。熟悉动力论的读者将在\eqref{align:ch06-6-7-13}中认出经典 Boltzmann 方程 BGK 逼近的模型。有望当 $\mu\downarrow 0$ 时，右端的刚性项将迫使 $f(v;x,t)$“松弛”到满足\eqref{align:ch06-6-7-2}的 $\chi(v;u(x,t))$。在验证这一预期能否实现之前，先讨论\eqref{align:ch06-6-7-13}、\eqref{align:ch06-6-7-12}解的性质。

\begin{theorem}\label{theorem:ch06-6-7-3}
设 $u_{0}\in L^{\infty}(\mathbb{R}^{m})\cap L^{1}(\mathbb{R}^{m})$。对任意 $\mu>0$，存在有界可测函数
\begin{align}
f_{\mu}(\cdot;\cdot,t)\in C^{0}([0,\infty);L^{1}(\mathbb{R}\times\mathbb{R}^{m})),\qquad
u_{\mu}(\cdot,t)\in C^{0}([0,\infty);L^{1}(\mathbb{R}^{m})),
\label{align:ch06-6-7-14}
\end{align}
它们是\eqref{align:ch06-6-7-13}、\eqref{align:ch06-6-7-12}以由 $u_{0}$ 诱导的初值
\begin{align}
f_{\mu}(v;x,0)=\chi(v;u_{0}(x)),\qquad v\in(-\infty,\infty),\ x\in\mathbb{R}^{m}
\label{align:ch06-6-7-15}
\end{align}
为初值的唯一解。对任意 $(x,t)\in\mathbb{R}^{m}\times[0,\infty)$，
\begin{align}
f_{\mu}(v;x,t)\in
\begin{cases}
[0,1], & v>0,\\
[-1,0], & v<0.
\end{cases}
\label{align:ch06-6-7-16}
\end{align}
若 $(\bar{f}_{\mu},\bar{u}_{\mu})$ 是\eqref{align:ch06-6-7-13}、\eqref{align:ch06-6-7-12}以 $\bar{u}_{0}\in L^{\infty}(\mathbb{R}^{m})\cap L^{1}(\mathbb{R}^{m})$ 诱导的初值的另一解，则对任意 $t>0$，
\begin{align}
\|f_{\mu}(\cdot;\cdot,t)-\bar{f}_{\mu}(\cdot;\cdot,t)\|_{L^{1}(\mathbb{R}\times\mathbb{R}^{m})}\leqslant\|f_{\mu}(\cdot;\cdot,0)-\bar{f}_{\mu}(\cdot;\cdot,0)\|_{L^{1}(\mathbb{R}\times\mathbb{R}^{m})},
\label{align:ch06-6-7-17}
\end{align}
\begin{align}
\|u_{\mu}(\cdot,t)-\bar{u}_{\mu}(\cdot,t)\|_{L^{1}(\mathbb{R}^{m})}\leqslant\|u_{0}(\cdot)-\bar{u}_{0}(\cdot)\|_{L^{1}(\mathbb{R}^{m})}.
\label{align:ch06-6-7-18}
\end{align}
此外，若
\begin{align}
u_{0}(x)\leqslant\bar{u}_{0}(x),\qquad x\in\mathbb{R}^{m},
\label{align:ch06-6-7-19}
\end{align}
则
\begin{align}
f_{\mu}(v;x,t)\leqslant\bar{f}_{\mu}(v;x,t),\qquad v\in(-\infty,\infty),\ x\in\mathbb{R}^{m},\ t\in[0,\infty),
\label{align:ch06-6-7-20}
\end{align}
\begin{align}
u_{\mu}(x,t)\leqslant\bar{u}_{\mu}(x,t),\qquad x\in\mathbb{R}^{m},\ t\in[0,\infty).
\label{align:ch06-6-7-21}
\end{align}
\end{theorem}

\begin{proof}
暂时承认 $(f_{\mu},u_{\mu})$ 与 $(\bar{f}_{\mu},\bar{u}_{\mu})$ 的存在性，沿特征线 $\mathrm{d}x/\mathrm{d}t=G^{\prime}(v)^{\top}$、$\mathrm{d}v/\mathrm{d}t=0$ 积分\eqref{align:ch06-6-7-13}，得
\begin{align}
f_{\mu}(v;x,t)=e^{-t/\mu}f_{\mu}(v;x-tG^{\prime}(v)^{\top},0)+\frac{1}{\mu}\int_{0}^{t}e^{-(t-\tau)/\mu}\chi(v;u_{\mu}(x-(t-\tau)G^{\prime}(v)^{\top},\tau))\,\mathrm{d}\tau.
\label{align:ch06-6-7-22}
\end{align}
于是由\eqref{align:ch06-6-7-22}、\eqref{align:ch06-6-7-15}及 $\chi$ 的性质，\eqref{align:ch06-6-7-16}直接成立。

对另一解 $(\bar{f}_{\mu},\bar{u}_{\mu})$ 写出\eqref{align:ch06-6-7-22}的类比式，从\eqref{align:ch06-6-7-22}中减去，得
\begin{align}
&f_{\mu}(v;x,t)-\bar{f}_{\mu}(v;x,t)=e^{-t/\mu}[f_{\mu}(v;x-tG^{\prime}(v)^{\top},0)-\bar{f}_{\mu}(v;x-tG^{\prime}(v)^{\top},0)]\notag\\
&\qquad+\frac{1}{\mu}\int_{0}^{t}e^{-(t-\tau)/\mu}[\chi(v;u_{\mu}(x-(t-\tau)G^{\prime}(v)^{\top},\tau))-\chi(v;\bar{u}_{\mu}(x-(t-\tau)G^{\prime}(v)^{\top},\tau))]\,\mathrm{d}\tau,
\label{align:ch06-6-7-23}
\end{align}
由此，利用
\begin{align}
\int_{-\infty}^{\infty}|\chi(v;u_{\mu})-\chi(v;\bar{u}_{\mu})|\,\mathrm{d}v\leqslant|u_{\mu}-\bar{u}_{\mu}|,
\label{align:ch06-6-7-24}
\end{align}
（它由\eqref{align:ch06-6-7-3}得到）并回顾\eqref{align:ch06-6-7-12}，得
\begin{align}
\|f_{\mu}(\cdot;\cdot,t)-\bar{f}_{\mu}(\cdot;\cdot,t)\|_{L^{1}(\mathbb{R}\times\mathbb{R}^{m})}
&\leqslant e^{-t/\mu}\|f_{\mu}(\cdot;\cdot,0)-\bar{f}_{\mu}(\cdot;\cdot,0)\|_{L^{1}(\mathbb{R}\times\mathbb{R}^{m})}\notag\\
&\quad+\frac{1}{\mu}\int_{0}^{t}e^{-(t-\tau)/\mu}\|\chi(v;u_{\mu}(x-(t-\tau)G^{\prime}(v)^{\top},\tau))\notag\\
&\qquad\qquad\qquad-\chi(v;\bar{u}_{\mu}(x-(t-\tau)G^{\prime}(v)^{\top},\tau))\|_{L^{1}(\mathbb{R}\times\mathbb{R}^{m})}\,\mathrm{d}\tau\notag\\
&\leqslant e^{-t/\mu}\|f_{\mu}(\cdot;\cdot,0)-\bar{f}_{\mu}(\cdot;\cdot,0)\|_{L^{1}(\mathbb{R}\times\mathbb{R}^{m})}\notag\\
&\quad+(1-e^{-t/\mu})\max_{0\leqslant\tau\leqslant t}\|f_{\mu}(\cdot;\cdot,\tau)-\bar{f}_{\mu}(\cdot;\cdot,\tau)\|_{L^{1}(\mathbb{R}\times\mathbb{R}^{m})}.
\label{align:ch06-6-7-25}
\end{align}
显然\eqref{align:ch06-6-7-25}蕴含\eqref{align:ch06-6-7-17}，后者进而给出\eqref{align:ch06-6-7-18}。特别地，\eqref{align:ch06-6-7-13}、\eqref{align:ch06-6-7-12}、\eqref{align:ch06-6-7-15}至多一解。此外，此解可由积分方程\eqref{align:ch06-6-7-22}通过 Picard 迭代构造。

由于 $\chi(v;u)$ 关于 $u$ 单调递增，\eqref{align:ch06-6-7-23}与\eqref{align:ch06-6-7-12}保证\eqref{align:ch06-6-7-19}蕴含\eqref{align:ch06-6-7-20}与\eqref{align:ch06-6-7-21}。证明完成。
\end{proof}

现在转向解当 $\mu\downarrow 0$ 时的极限行为。

\begin{theorem}\label{theorem:ch06-6-7-4}
对 $\mu>0$，设 $(f_{\mu},u_{\mu})$ 为\eqref{align:ch06-6-7-13}、\eqref{align:ch06-6-7-12}、\eqref{align:ch06-6-7-15}以 $u_{0}\in L^{\infty}(\mathbb{R}^{m})\cap L^{1}(\mathbb{R}^{m})$ 为初值的解。则当 $\mu\downarrow 0$ 时，
\begin{align}
u_{\mu}(x,t)\to u(x,t),
\label{align:ch06-6-7-26}
\end{align}
\begin{align}
f_{\mu}(v;x,t)\to\chi(v;u(x,t)),
\label{align:ch06-6-7-27}
\end{align}
在 $L^{1}_{\mathrm{loc}}$ 中，其中 $\chi(v;u)$ 对某有界非负测度 $\nu$ 满足\eqref{align:ch06-6-7-2}，因而 $u$ 是\eqref{align:ch06-6-1-1}、\eqref{align:ch06-6-1-2}的可容许弱解。
\end{theorem}

\begin{proof}
第一步是证明族 $\{(f_{\mu},u_{\mu}):\mu>0\}$ 在平均意义下等度连续。由压缩性质\eqref{align:ch06-6-7-17}、\eqref{align:ch06-6-7-18}，在 $v$ 与 $x$ 方向显然如此。对任意 $w\in\mathbb{R}$、$y\in\mathbb{R}^{m}$，由 $\bar{f}_{\mu}(v;x,t)=f_{\mu}(v+w;x+y,t)$、$\bar{u}_{\mu}(x,t)=u_{\mu}(x+y,t)$ 定义的函数 $(\bar{f}_{\mu},\bar{u}_{\mu})$ 是\eqref{align:ch06-6-7-13}、\eqref{align:ch06-6-7-12}以 $\bar{f}_{\mu}(v;x,0)=\chi(v+w;u_{0}(x+y))$ 为初值的解，故
\begin{align}
\int_{0}^{\infty}\int_{\mathbb{R}^{m}}|f_{\mu}(v+w;x+y,t)-f_{\mu}(v;x,t)|\,\mathrm{d}x\,\mathrm{d}v
\leqslant\int_{0}^{\infty}\int_{\mathbb{R}^{m}}|\chi(v+w;u_{0}(x+y))-\chi(v;u_{0}(x))|\,\mathrm{d}x\,\mathrm{d}v,
\label{align:ch06-6-7-28}
\end{align}
\begin{align}
\int_{\mathbb{R}^{m}}|u_{\mu}(x+y,t)-u_{\mu}(x,t)|\,\mathrm{d}x\leqslant\int_{\mathbb{R}^{m}}|u_{0}(x+y)-u_{0}(x)|\,\mathrm{d}x.
\label{align:ch06-6-7-29}
\end{align}
$t$ 方向的等度连续由上述结果结合输运方程\eqref{align:ch06-6-7-13}本身易得；细节从略。

接着考虑函数
\begin{align}
\omega_{\mu}(v;x,t)=\int_{-\infty}^{v}[\chi(w;u_{\mu}(x,t))-f_{\mu}(w;x,t)]\,\mathrm{d}w.
\label{align:ch06-6-7-30}
\end{align}
固定 $(x,t)$，为确定起见设 $u_{\mu}(x,t)>0$（其他情形类似处理）。显然 $\omega_{\mu}(-\infty;x,t)=0$。由\eqref{align:ch06-6-7-3}与\eqref{align:ch06-6-7-16}，$\omega_{\mu}(\cdot;x,t)$ 在区间 $(-\infty,u_{\mu}(x,t))$ 上非减，在区间 $(u_{\mu}(x,t),\infty)$ 上非增。最后由\eqref{align:ch06-6-7-12}，$\omega_{\mu}(\infty;x,t)=0$。因而可写
\begin{align}
\frac{1}{\mu}[\chi(v;u_{\mu}(x,t))-f_{\mu}(v;x,t)]=\partial_{v}\nu_{\mu},
\label{align:ch06-6-7-31}
\end{align}
其中 $\nu_{\mu}$ 是非负测度，关于 $\mu>0$ 一致有界。

于是从任一序列 $\{\mu_{k}\}$（$\mu_{k}\to 0$，$k\to\infty$）可提取子列（仍记为 $\{\mu_{k}\}$），使 $\{(f_{\mu_{k}},u_{\mu_{k}})\}$ 在 $L^{1}_{\mathrm{loc}}$ 中收敛于函数 $(f,u)$，且 $\{\nu_{\mu_{k}}\}$ 在测度空间中弱收敛于有界非负测度 $\nu$。显然 $f(v;x,t)=\chi(v;u(x,t))$，且\eqref{align:ch06-6-7-2}成立。由唯一性，整个族 $\{(f_{\mu},u_{\mu})\}$ 当 $\mu\downarrow 0$ 时收敛于 $(\chi(\cdot;u),u)$。证明完成。
\end{proof}

有趣的是，输运方程\eqref{align:ch06-6-1-9}（正如在\cref{section:ch06-cauchy问题经典解}中所见，它生成\eqref{align:ch06-6-1-1}的经典解）也以\eqref{align:ch06-6-7-13}或\eqref{align:ch06-6-7-2}的面貌出现在弱解理论中。

值得注意的是，\eqref{align:ch06-6-1-9}与\eqref{align:ch06-6-1-1}的弱解之间还有另一联系，它由下述考虑引出。设 $u(x,t)$ 是\eqref{align:ch06-6-1-1}、\eqref{align:ch06-6-1-2}的可容许弱解，初值 $u_{0}$ 关于 $x_{\alpha}$ 为 $1$-周期，取值于区间 $[a,b]\subset(0,1)$。于是 $u(\cdot,t)$ 关于 $x_{\alpha}$ 也为 $1$-周期，取值于 $[a,b]$。记 $\mathscr{C}$ 为 $\mathbb{R}^{m}$ 中的闭单位立方体 $\mathscr{C}=[0,1]^{m}$。目的是通过 $u(\cdot,t)$ 在 $\mathscr{C}$ 中的等值集来刻画它。为此，引入定义于 $v\in[0,1]$、$x\in\mathscr{C}$、$t\geqslant 0$ 的函数 $f(v;x,t)$，规定
\begin{align}
\{x\in\mathscr{C}:f(v;x,t)\leqslant 1\}=\{x\in\mathscr{C}:u(x,t)\geqslant v\}.
\label{align:ch06-6-7-32}
\end{align}
可由 $f$ 经
\begin{align}
u(x,t)=\int_{0}^{1}h(1-f(v;x,t))\,\mathrm{d}v
\label{align:ch06-6-7-33}
\end{align}
恢复 $u$，其中 $h$ 表示 Heaviside 函数，$y<0$ 时 $h(y)=0$，$y\geqslant 0$ 时 $h(y)=1$。

正如\cref{section:ch06-注记}所引文献所示，$f$ 满足 Hilbert 空间 $H=L^{2}([0,1]\times\mathscr{C})$ 上的抽象微分方程

\begin{align}
\partial_{t}f+\sum_{\alpha=1}^{m}Q_{\alpha}^{\prime}(v)\partial_{\alpha}f+\partial\mathscr{K}(f)\ni 0,
\label{align:ch06-6-7-34}
\end{align}
其中 $\partial\mathscr{K}$ 表示闭凸锥 $\mathscr{K}=\{f\in H:f_{v}\geqslant 0\}$ 的次微分，即
\begin{align}
\partial\mathscr{K}(f)=\Bigl{\{}g\in H:\int_{0}^{1}\int_{\mathscr{C}}(\bar{f}-f)g\,\mathrm{d}x\,\mathrm{d}v\leqslant 0,\ \text{对一切 }\bar{f}\in\mathscr{K}\Bigr{\}}.
\label{align:ch06-6-7-35}
\end{align}
上述观察的用处在于，方程\eqref{align:ch06-6-7-34}涉及极大单调算子，因而在 $H$ 上生成压缩半群。易见可调整初值 $f(v;x,0)$，使经\eqref{align:ch06-6-7-33}计算出的 $u(x,0)$ 与给定初值 $u_{0}(x)$ 重合。因而 $f$（以及由此 $u$）的存在性、唯一性、稳定性乃至数值构造，都可由 Hilbert 空间中压缩半群的标准泛函分析理论得到。此途径的细节及其与动力论表述的联系，见\cref{section:ch06-注记}所引文献。


\end{document}
